\chapter{Ova and Parasite (O\&P) Examination}

\section*{What Is This Test?}
The ova and parasite (O\&P) examination is a microscopic evaluation of stool for the detection and identification of intestinal parasites, including protozoa (unicellular: trophozoites and cysts) and helminths (multicellular worms: ova/eggs, larvae, adult worms, and proglottids). Protozoa causing GI disease include \textit{Giardia duodenalis} (\textit{G. lamblia}), \textit{Entamoeba histolytica}, \textit{Cryptosporidium} species, \textit{Cystoisospora belli}, \textit{Cyclospora cayetanensis}, \textit{Dientamoeba fragilis}, and \textit{Balantidium coli}. Helminths include nematodes (\textit{Ascaris lumbricoides}, \textit{Trichuris trichiura}, hookworm [\textit{Ancylostoma duodenale}, \textit{Necator americanus}], \textit{Strongyloides stercoralis}, \textit{Enterobius vermicularis}), cestodes (\textit{Taenia saginata}, \textit{T. solium}, \textit{Diphyllobothrium latum}, \textit{Hymenolepis nana}), and trematodes (\textit{Schistosoma} species, \textit{Fasciola hepatica}, \textit{Clonorchis sinensis}, \textit{Paragonimus westermani}). The O\&P uses macroscopic and microscopic analysis of concentrated and permanently stained preparations \cite{garcia2016diagnostic}. Multiple specimens (3 stools on alternate days over 10 days) are recommended because parasite shedding is intermittent. Sensitivity of a single specimen is 50--70\%; 3 specimens increase sensitivity to 85--95\%.

\section*{Alternative Names}
Stool O\&P; Parasitology Examination; Fecal Parasite Exam; Stool for Parasites; Ova and Parasite Screen; Stool for Ova; O\&P Concentrate and Trichrome Stain

\section*{Principle}
The O\&P examination comprises: (1) Macroscopic examination: Visual inspection for consistency, blood, mucus, and macroscopic parasites (adult \textit{Ascaris} worms, tapeworm proglottids). (2) Direct wet mount: Fresh stool mixed with saline (for motile trophozoites---ameboid \textit{Entamoeba}, falling-leaf \textit{Giardia}) and iodine (for nuclear detail of cysts). (3) Concentration technique: Formalin-ethyl acetate (FEA) sedimentation concentrates helminth ova, larvae, and protozoan cysts from 2--3 grams of stool \cite{cdc2024dpdx}. Zinc sulfate flotation (specific gravity 1.18--1.20) floats lighter cysts to the surface. (4) Permanently stained smear: Trichrome stain (Wheatley modification) is the gold standard for morphologic identification of intestinal protozoa. Nuclei, chromatoid bodies, and RBCs stain red/pink; cytoplasm stains blue-green to purple. The entire 22 x 22 mm coverslip is examined at 1,000x. (5) Special stains: Modified acid-fast (Kinyoun) for \textit{Cryptosporidium} (4--6 \textmu m, red/pink), \textit{Cystoisospora} (20--30 \textmu m), \textit{Cyclospora} (8--10 \textmu m). Modified trichrome (Weber) for microsporidia (1--2 \textmu m). Cellophane tape test for \textit{Enterobius vermicularis} (pinworm) eggs (50--60 \textmu m, D-shaped).

\section*{History}
Intestinal parasites have been recognized since antiquity. Van Leeuwenhoek identified \textit{Giardia} in his own stool in 1681. The formalin-ether (now formalin-ethyl acetate) concentration method was developed in the early 20th century. Trichrome staining was developed in the 1950s. Antigen detection tests (EIA) for \textit{Giardia} and \textit{Cryptosporidium} were introduced in the 1990s and are more sensitive than O\&P for these pathogens. Multiplex PCR panels (FilmArray GI, Luminex xTAG GPP) were introduced in the 2010s and detect protozoan pathogens with higher sensitivity, but commercial panels do not detect helminths \cite{mchardy2014detection}. The O\&P remains indispensable for detection of helminths and for identification of parasites not included in molecular panels.

\section*{Why This Test Is Ordered}
\begin{itemize}
  \item Diagnosis of intestinal parasitic infections in patients with persistent diarrhea (>2 weeks), dysentery, abdominal pain, bloating, nausea, weight loss, or unexplained eosinophilia
  \item Investigation of GI symptoms in travelers returning from tropical countries, immigrants, and refugees
  \item Evaluation of immunocompromised patients (HIV/AIDS, transplant, hematologic malignancy) with diarrhea; \textit{Cryptosporidium}, \textit{Cystoisospora}, \textit{Cyclospora}, microsporidia, and \textit{Strongyloides} hyperinfection syndrome are particularly prevalent and severe
  \item Evaluation of unexplained eosinophilia: tissue-invasive helminths (\textit{Strongyloides}, hookworm, \textit{Ascaris} larval migration, \textit{Trichinella}, \textit{Toxocara}) induce eosinophilia, while luminal protozoa generally do not
  \item Screening asymptomatic high-risk individuals (refugees, internationally adopted children) for chronic parasitic infections
  \item Investigation of perianal pruritus, particularly nocturnal, for \textit{Enterobius vermicularis} (pinworm)
  \item Evaluation of suspected \textit{Strongyloides stercoralis} infection in patients from endemic areas prior to immunosuppression (serology is preferred for screening; stool O\&P has low sensitivity ~30--50\% even with multiple specimens)
\end{itemize}

\section*{Primary vs Secondary Test}
The O\&P examination is the primary traditional diagnostic test for intestinal parasitic infections. For protozoan pathogens (\textit{Giardia}, \textit{Cryptosporidium}, \textit{Entamoeba histolytica}), antigen detection (EIA) and multiplex PCR panels are now preferred as primary tests due to higher sensitivity and faster turnaround. The O\&P remains the primary test for detection of helminths and for species-level identification not provided by molecular panels. For \textit{Strongyloides}, serology (Strongyloides IgG ELISA) is the preferred screening test.

\section*{How This Test Is Performed}
Collection of 3 stool specimens on alternate days (within a 10-day period) is recommended. Fresh stool should be collected directly into a clean, dry, leak-proof container (not contaminated with urine or water). The specimen should be transported to the laboratory within 1--2 hours of passage if fresh examination for trophozoites is requested (trophozoites lose motility rapidly). Alternatively, stool can be placed into a two-vial collection kit: one vial with 5--10\% formalin for concentration and wet mounts; the other with PVA (polyvinyl alcohol) fixative for permanent trichrome-stained smears. Specimens in fixatives are stable for months at room temperature \cite{garcia2018practical}. The Scotch tape test for pinworm: applied to perianal skin early in the morning on 3 consecutive mornings before bathing or defecation; the tape is applied sticky-side down to a glass slide.

\section*{What Preparation Is Needed}
Patients should avoid barium and bismuth-containing compounds for 5--7 days, and mineral oil and non-absorbable antidiarrheal preparations for 5--7 days before specimen collection. Antibiotics can temporarily suppress some parasites and should be avoided for 2--3 weeks before screening, though testing should not be delayed in acutely ill patients. For the Scotch tape test, the specimen should be collected in the morning before bathing or using the toilet.

\section*{Contraindications and Precautions}
The O\&P examination is labor-intensive and requires extensive expertise. Sensitivity depends heavily on the microscopist's skill. A negative O\&P does NOT definitively exclude parasitic infection. Parasite shedding is intermittent; even with 3 specimens, negative predictive value is ~85--95\%. \textit{Entamoeba histolytica} (pathogenic, causes amebic dysentery and liver abscess) is morphologically identical to \textit{E. dispar} (nonpathogenic, approximately 10 times more common) on O\&P. The finding of "\textit{Entamoeba histolytica/dispar}" requires confirmatory \textit{E. histolytica}-specific antigen detection or PCR if the patient is symptomatic. Erythrophagocytosis (ingested RBCs within trophozoites) strongly suggests \textit{E. histolytica}. \textit{Blastocystis hominis} pathogenicity remains controversial. Fecal leukocytes (PMNs) indicate inflammation and suggest bacterial dysentery or \textit{E. histolytica}. Macrophages may be mistaken for \textit{Entamoeba} trophozoites by inexperienced microscopists.

\section*{Risks}
No risks from stool collection. The primary clinical risk is failure to diagnose \textit{Strongyloides stercoralis} before initiating immunosuppression (corticosteroids, transplant, anti-TNF therapy). Disseminated strongyloidiasis/hyperinfection syndrome has mortality of 60--85\%. Given the low sensitivity of stool O\&P for \textit{Strongyloides}, serology should be obtained in at-risk patients prior to planned immunosuppression.

\section*{What Happens After the Test}
Macroscopic and direct wet mount results available same day. Concentrated wet mounts: 24 hours. Trichrome-stained smear: 24--48 hours. Special stains (modified AFB, modified trichrome): additional 24--48 hours. Final report: typically 48--72 hours. All detected parasites are reported, including nonpathogenic protozoa (\textit{E. dispar}, \textit{E. coli}, \textit{Endolimax nana}, \textit{Iodamoeba butschlii}, \textit{Chilomastix mesnili}) that indicate fecal-oral exposure but do not require treatment.

\section*{Understanding Results}

\subsection*{Below Normal}
\begin{itemize}
  \item \textbf{No ova or parasites seen:} No parasites detected. With 3 specimens, NPV reaches 85--95\%, but a negative O\&P does NOT absolutely exclude parasitic infection. If suspicion persists, consider serology (\textit{Strongyloides} IgG), antigen detection (\textit{Giardia}, \textit{Cryptosporidium}), multiplex PCR GI panel, or specialty tests (Scotch tape test, Baermann concentration, agar plate culture, duodenal aspirate)
  \item \textbf{Nonpathogenic protozoa detected (\textit{E. dispar}, \textit{E. coli}, \textit{Endolimax nana}, etc.):} Indicates fecal-oral exposure but does not cause disease. No treatment indicated. However, their presence signals conditions for transmission of pathogenic organisms; consider testing for co-infecting pathogens
\end{itemize}

\subsection*{Above Normal}
\begin{itemize}
  \item \textbf{\textit{Giardia duodenalis}:} Pear-shaped trophozoites with 2 anterior nuclei and ventral sucking disk; oval cysts (8--12 \textmu m) with 2--4 nuclei. Diagnostic of giardiasis. Treatment: tinidazole 2 g single dose, or metronidazole 500 mg BID x 5--7 days, or nitazoxanide 500 mg BID x 3 days
  \item \textbf{\textit{Entamoeba histolytica/dispar}:} Cysts with 1--4 nuclei; trophozoites with "ring-and-dot" nucleus. Must be differentiated as \textit{E. histolytica} (treat) vs. \textit{E. dispar} (no treatment needed) by antigen or PCR. If confirmed \textit{E. histolytica}, treat with metronidazole/tinidazole followed by a luminal agent (paromomycin or iodoquinol)
  \item \textbf{\textit{Cryptosporidium} species} (4--6 \textmu m oocysts, red on modified AFB): Common in HIV with CD4 <200. Nitazoxanide for immunocompetent; immune reconstitution with ART is cornerstone in HIV \cite{checkley2015review}
  \item \textbf{Helminth ova:} \textit{Ascaris lumbricoides} (oval, 45--70 \textmu m, mammillated shell), \textit{Trichuris trichiura} (barrel-shaped, 50--55 \textmu m, bipolar mucus plugs), hookworm (thin-shelled, oval, 60--75 \textmu m), \textit{Taenia} species (spherical, 30--40 \textmu m, radially striated shell---cannot distinguish \textit{T. solium} from \textit{T. saginata} by egg alone), \textit{Schistosoma} species (large, with characteristic spine---terminal in \textit{S. haematobium}, lateral in \textit{S. mansoni}). Treatment is species-specific (albendazole, praziquantel, ivermectin)
  \item \textbf{\textit{Strongyloides stercoralis} larvae (rhabditiform, 250--300 \textmu m):} Diagnostic. Treatment: ivermectin 200 mcg/kg/day x 2 days (uncomplicated). Disseminated disease requires prolonged ivermectin and possible rectal/veterinary formulations
\end{itemize}

\section*{Normal Results}
\textbf{No ova or parasites seen.} No pathogenic parasites detected in the concentrated and permanently stained preparations after examination of 3 specimens collected on alternate days.

\section*{Follow-up Tests}
\begin{itemize}
  \item \textbf{\textit{Giardia} and \textit{Cryptosporidium} stool antigen EIA or DFA:} More sensitive than O\&P; should be ordered when O\&P is negative but clinical suspicion persists for these pathogens
  \item \textbf{Multiplex GI pathogen PCR panel (FilmArray GI, Luminex xTAG GPP):} Detects protozoan pathogens with higher sensitivity; limitation: does not detect helminths
  \item \textbf{\textit{Strongyloides} IgG serology:} Preferred screening test for strongyloidiasis (high sensitivity >90\%); stool O\&P is insensitive
  \item \textbf{Scotch tape test for \textit{Enterobius vermicularis}:} 3 consecutive morning specimens; sensitivity 50--90\%. Empiric treatment (albendazole or mebendazole) is often given without testing
  \item \textbf{Baermann funnel technique or agar plate culture:} For \textit{Strongyloides} when stool O\&P is negative but serology is positive; higher sensitivity than routine concentration
  \item \textbf{CBC with differential:} Eosinophilia (>500 cells/$\mu$L) supports tissue-invasive helminth infections; not typically present in luminal protozoan infections
\end{itemize}

\section*{References}
\bibliographystyle{unsrtnat}
\bibliography{references}

\clearpage
\section*{Regulatory Status and Authorization Date}
Regulatory status applies to the specific assay, instrument, manufacturer, and intended use rather than to the test name alone. No single approval date can be assigned to this test category. For a commercial assay, verify the current product in the relevant regulator's database: the U.S. Food and Drug Administration (FDA) clearance, approval, or authorization; India's Central Drugs Standard Control Organisation (CDSCO) licence or permission; the European Union In Vitro Diagnostic Regulation (EU IVDR) CE certificate issued through the applicable conformity-assessment route; or the United Kingdom Medicines and Healthcare products Regulatory Agency (MHRA) registration and UKCA/CE status. Laboratory-developed tests may not have an individual FDA, CDSCO, EU IVDR, or MHRA approval date.
\clearpage
